\documentclass[11pt,a4paper]{article}
\usepackage[T1]{fontenc}
\usepackage{amsmath,amssymb,geometry,booktabs}
\geometry{margin=27mm}
\title{EPPA--III: Section 3 verification checkpoint}
\author{Independent formalization audit; source manuscript frozen}
\date{12 October 2026}
\begin{document}
\maketitle

\section*{Purpose and verification boundary}

This checkpoint concerns the direct alternative proof of
Lemma~3.7 for a double-sized EPPA witness of
\(G=sK_t\), with the exceptional Wagner graph when
\(s=t=2\).  The original manuscript is frozen:
all proposed corrections remain in the separate validation files.
A successful Lean build and the repository's explicit permitted-axiom
audit are required before a new theorem is called formally verified.

\section*{Verified results merged into \texttt{main}}

The following unrestricted mathematical statements have completed
Lean CI and were merged (the PR numbers identify the audit trail).

\begin{itemize}
\item The full fixed-coordinate-size wreath-orbit bound and its
  inclusion in the actual EPPA neighbourhood-profile orbit;
  extreme partial sizes \(k=1,t-1\)
  (PRs~\#29, \#32 and earlier prerequisites).
\item Repeated partial-coordinate degrees force \(s=t=2\);
  independent inner permutations and saturation exclude
  mixed partial sizes; thus \emph{any two proper partial
  coordinates imply \(s=t=2\)}
  (PRs~\#35, \#38, \#46).
\item Saturation under outer clique permutations forces the
  unselected coordinate slices to agree; for a unique proper
  partial coordinate those slices are uniformly empty or full
  (PRs~\#47 and \#49).
\item In the zero-partial case the selected cliques consist of
  true twins and occupy pairwise distinct classes of equality
  of closed neighbourhoods (PRs~\#50, \#54).
  Host automorphisms carry these classes to one another,
  so a transitive host has equal class sizes
  (PR~\#56).
\item The canonical true-twin quotient is a genuine loopless
  simple graph with well-defined adjacency and descending
  automorphisms. The selected clique classes form an
  independent set of cardinal \(s\). All these quotient
  interfaces, together with equality of the two induced
  half-edge counts in a vertex-transitive labelled host,
  were integrated in PR~\#68.
\item A positive-degree \(d\)-regular graph has no independent
  set greater than half its order (PR~\#62).
  If the host has \(2st\) vertices and the quotient contains an edge,
  quotient transitivity and double-counting then force
  exactly \(2s\) twin classes of order \(t\)
  (PR~\#71).
\item If a spanning subgraph and its supergraph have equal
  numbers of edges, they coincide. Its specialization to
  \(sK_t\) is verified in PR~\#74.
\end{itemize}

These theorems do \emph{not} yet constitute a complete
formal proof of Lemma~3.7. In particular the last graph
classification, which passes from twin classes to the
uniform connection matrix, remains a separate obligation.

\section*{Current source branches}

As of this checkpoint, the outstanding Lean checks include:
\begin{itemize}
\item PR~\#75: a non-cluster graph has an edge in its canonical
  twin quotient, removing the positive-degree hypothesis
  from the zero-partial counting theorem;
  PR~\#76 composes this with PR~\#71.
\item PR~\#73: a selected \(K_t\) exhausts any containing
  twin class of size \(t\); PR~\#78 then identifies
  the outside half with the unselected quotient fibers.
\item PR~\#77: from \(\binom sr\le s\) deduce
  \(r\in\{0,1,s-1,s\}\), uniformly in the parameters.
\item PR~\#79: when a profile choice orbit saturates the
  \(st\)-vertex budget, all outside vertices have distinct
  profiles and the outside-profile image is precisely
  that orbit. Its unique-proper-coordinate specialization
  prepares the natural wreath action on the outside half.
\end{itemize}
The status of these PRs must be checked against CI before merging;
an earlier green result does not certify a later rewritten branch.

\section*{Remaining structural assembly}

\paragraph{Zero partial coordinates.}
Once the quotient has \(2s\) classes of size \(t\),
each selected \(K_t\) fills one entire class.
The remaining \(s\) classes partition the complementary half \(B\)
and each induces \(K_t\). Because regularity gives
\[
    |E(H[B])|=|E(G)|=s\binom t2 ,
\]
there are no edges between these latter classes, hence
\(H[B]\cong sK_t\).
Every \(B\)--\(G\) connection block is complete or empty,
as its endpoints lie in twin classes.
Vertex-transitivity makes the number \(r\) of selected
cliques met by a vertex of \(B\) independent of that vertex.
The base-clique permutation subgroup, whose automorphisms
extend by EPPA, produces all
\(\binom sr\) masks; there are only \(s\) complementary
twin classes. Thus \(\binom sr\le s\) and
\(r\in\{0,1,s-1,s\}\).
For \(r=1,s-1\), all singleton or co-singleton masks
occur exactly once, producing the uniform diagonal/off-diagonal
connection matrix. For \(r=0,s\), all cross entries are uniform.
The classification of the graph from these observations
is established on paper but the final Lean assembly is open.

The assumption that the quotient has an edge is automatic if
\(H\) is not a cluster graph. If \(H\) \emph{is} a cluster graph,
vertex-transitivity makes every connected clique component
equinumerous; hence \(H=mK_a\), a homogeneous graph.
This supplies a short direct paper justification that a
non-homogeneous vertex-transitive host is not a cluster graph,
without invoking the full finite homogeneous graph classification.
The corresponding homogeneity extension argument is not yet Lean-checked.

\paragraph{One proper partial coordinate.}
The fixed-size profile orbit has exactly \(st\) possible
choices, and the unique outside profile corresponding to each
choice must be established as a finite bijection (PR~\#79).
The induced canonical action is the imprimitive
\(S_t\wr S_s\)-action on \(s\) blocks of size \(t\);
its unordered pair-orbits are the same-block and
distinct-block pairs. By the equal-edge-count identity,
\(H[B]\cong sK_t\), giving matching or co-matching diagonal
connections and uniform off-diagonal connections.
Formalization of the actual induced action and its pair-orbits
is the key remaining step in this case.

\paragraph{Wagner exception.}
When \(s=t=2\) and two coordinates are partial,
the four outside profiles form the corners of a square.
The induced wreath group has two unordered-pair orbits:
four sides and two opposite-corner pairs.
The equal-edge count is two, so only the opposite
matching can occur. It yields the Wagner graph and the
connection matrix
\[
  \begin{pmatrix} M&M\\ M&\overline M\end{pmatrix}.
\]
The graph-level identification still requires a short Lean
orbit-count argument; the earlier cube obstruction is
already verified in \texttt{WagnerReduction.lean}.

\paragraph{Low parameters and homogeneous hosts.}
The \(s=1\) and \(t=1\) reductions, and the homogeneous
exceptional cases, remain at paper level in
\texttt{validation/direct-clique-cover-proof.tex}.
They should not be included in the claim of complete
Lean verification until independently certified.

\section*{Editorial policy}
Do not modify the author manuscript until a correction is
reviewed. Preserve the short Wagner repair as the default
localized suggestion; the alternative direct proof is
recorded only as an independent validation route.
For every future PR, keep the validated theorem statement,
the hypotheses, its axiom report, and remaining dependencies
clearly separated.

\end{document}
